\chapter{Pricing e revenue management}
\label{cap:pricing}

\noindent\textbf{Classe:} NLP concavo / non convesso \hfill
\textbf{Script:} \texttt{python/lab07\_pricing.py}

\section{Motivazione gestionale}

Quale prezzo massimizza il profitto di un concerto, di un hotel, di un abbonamento,
quando la domanda diminuisce al crescere del prezzo? Qui il prezzo non è un dato ma
una \emph{variabile}: la domanda diventa endogena e il profitto $p \cdot q$ introduce
un termine bilineare. È il capitolo in cui il laboratorio incontra per la prima volta
la non convessità --- e impara a maneggiarla.



Fissare il prezzo è forse la decisione aziendale a più alto impatto sul margine: un
punto percentuale di prezzo vale spesso più di molti punti di volume. Il revenue
management --- nato nelle compagnie aeree e oggi standard in hotel, eventi,
noleggi --- aggiunge il vincolo di capacità: i posti invenduti stasera non si
recuperano domani.

\textbf{In questo capitolo impareremo a:} (1) modellare la domanda come funzione
del prezzo e capire quando il profitto è concavo; (2) trattare un obiettivo
bilineare non convesso; (3) valutare correttamente la capacità (il posto in più
\emph{non} vale il margine pieno); (4) riconoscere gli ottimi di spigolo, dove
nessuna derivata si annulla.

\section{Costruzione guidata del modello}

\paragraph{Il problema a parole.} \emph{Decidiamo} il prezzo di vendita e la
quantità da vendere. \emph{L'obiettivo} è massimizzare il profitto, cioè ricavi meno
costi variabili. \emph{I vincoli}: non si può vendere più di quanto il mercato chiede
a quel prezzo, né più della capacità disponibile.

\subsection*{Dati (input del modello)}
\begin{center}\small
\begin{tabular}{llp{7.4cm}}
\toprule
Simbolo & Tipo & Significato \\
\midrule
$d(\cdot)$ & funzione $\R_{\ge 0} \to \R_{\ge 0}$ & domanda attesa al prezzo $p$, decrescente \\
$c$ & $\in \R_{\ge 0}$ & costo marginale unitario \\
$k$ & $\in \R_{> 0}$ & capacità disponibile (posti, camere, unità) \\
\bottomrule
\end{tabular}
\end{center}

\subsection*{Variabili decisionali}
\begin{center}\small
\begin{tabular}{llp{7.4cm}}
\toprule
$p$ & $\ge 0$, continua & prezzo deciso dall'impresa \\
$q$ & $\ge 0$, continua & quantità venduta \\
\bottomrule
\end{tabular}
\end{center}

\begin{modello}[Pricing con vincolo di capacità]
\begin{align}
\max\;& p\,q - c\,q \notag\\
\text{s.t.}\;\; & q \le d(p) \label{eq:dom7}\\
& q \le k, \qquad p, q \ge 0 \notag
\end{align}
Funzioni di domanda usate nel caso di studio (i coefficienti sono tutti dati reali
positivi): lineare $d(p) = a - bp$; a elasticità costante
$d(p) = \theta\,p^{-\varepsilon}$ con $\varepsilon > 1$; logistica
$d(p) = m / \bigl(1 + e^{-\alpha + \beta p}\bigr)$.
\end{modello}

\begin{spiegazione}
\textbf{Funzione obiettivo.} Il profitto $p\,q - c\,q = (p - c)\,q$: ricavo meno
costo variabile. Contiene il prodotto tra due variabili ($p$ e $q$), ed è qui che
nasce tutta la difficoltà del capitolo.

\textbf{Vincolo di capacità ($q \le k$).} Non si può vendere più dei posti
disponibili; quando è attivo, è lui a determinare il prezzo ottimo.

\textbf{Perché $q \le d(p)$ e non $q = d(p)$?} All'ottimo la disuguaglianza
\eqref{eq:dom7} è attiva da sola (vendere meno del possibile a quel prezzo non
conviene mai se $p > c$), ma scriverla come $\le$ mantiene la regione ammissibile
più semplice e il modello ammissibile anche quando $d(p) > k$.

\textbf{Dove sta la non convessità?} Nel prodotto $p \cdot q$ dell'obiettivo: è un
termine bilineare, indefinito. Con domanda lineare, sostituendo $q = a - bp$, il
profitto $(p-c)(a-bp)$ è una parabola concava: il problema \emph{ridotto} è facile.
Gurobi risolve comunque la versione bilineare all'ottimo globale
(\texttt{NonConvex=2}); per domande generali usiamo \texttt{scipy} e uno studio di
funzione.
\end{spiegazione}

\section{Esempio svolto a mano}

\begin{esempio}[Domanda lineare, con e senza capacità]
Concerto: $d(p) = 1200 - 5p$ (quindi $a = 1200$, $b = 5$), costo marginale $c = 20$~\euro, teatro da $k = 400$ posti.

\textbf{Passo 1 — senza capacità.} Profitto
$\Pi(p) = (p - 20)(1200 - 5p)$; $\Pi'(p) = 1200 - 10p + 100 = 0$
dà $p^\circ = \dfrac{a/b + c}{2} = \dfrac{240 + 20}{2} = 130$~\euro,
con $q^\circ = 1200 - 650 = 550$ biglietti.

\textbf{Passo 2 — la capacità morde.} $q^\circ = 550 > 400$: il vincolo è attivo,
quindi si vende esattamente $k$ e il prezzo che svuota il teatro è
$p^* = (a - k)/b = (1200 - 400)/5 = 160$~\euro. Profitto:
$(160 - 20)\cdot 400 = 56.000$~\euro.

\textbf{Passo 3 — quanto vale un posto in più?} Aumentando $k$, il prezzo ottimo
scende di $1/b$: $\dfrac{d\Pi}{dk} = \underbrace{(p^* - c)}_{140}
+ k \underbrace{\dfrac{dp^*}{dk}}_{-1/5} = 140 - 80 = 60~\text{\euro}.$
Un posto in più \emph{non} vale il margine pieno (140~\euro): per riempirlo bisogna
abbassare il prezzo a tutti. Errore classico di valutazione degli investimenti di
capacità.
\end{esempio}

\section{Caso di studio e implementazione}

Stessi dati dell'esempio (è il caso raro in cui l'istanza realistica è già piccola),
più le varianti a elasticità costante ($\theta = 6 \cdot 10^6$, $\varepsilon = 2{,}2$),
logistica ($m = 900$, $\alpha = 6$, $\beta = 0{,}045$) e il multiprodotto.

\begin{lstlisting}[style=python]
m = gp.Model("pricing_lineare")
m.Params.NonConvex = 2            # obiettivo bilineare p*q
p = m.addVar(ub=a/b, name="p")
q = m.addVar(name="q")
m.addConstr(q <= a - b*p, name="domanda")
v_cap = m.addConstr(q <= K, name="capienza")
m.setObjective(p*q - c*q, GRB.MAXIMIZE)
m.optimize()
# niente duali nei QP non convessi: sensitivita' per perturbazione
v_cap.RHS = K + 1; m.optimize()
\end{lstlisting}

\begin{lstlisting}[style=output]
Gurobi:  p* = 160,00 EUR, q* = 400, profitto = 56.000,00 EUR
Valore marginale di un posto: 59,80 EUR (teoria: 60,00)
Elasticita' costante: p* = 79,11 EUR  (ottimo nel punto in cui d(p) = k)
Logistica           : p* = 138,29 EUR, profitto = 47.316,83 EUR
\end{lstlisting}

\begin{center}
\begin{tikzpicture}
\begin{axis}[labmezza, title={Profitto concavo e vincolo di capienza},
  xlabel={prezzo (\euro)}, ylabel={profitto (\euro)},
  legend style={at={(0.03,0.03)}, anchor=south west}]
\addplot [grigiomedio, dashed, thick] table [col sep=comma, x=p, y=senza]
  {\figdat{cap07_profitto}};
\addplot [teal, very thick] table [col sep=comma, x=p, y=con]
  {\figdat{cap07_profitto}};
\addplot [rossomattone, dash dot] coordinates {(160, 0) (160, 60500)};
\legend{senza capienza, con $k = 400$, $p^* = 160$}
\end{axis}
\end{tikzpicture}%
\begin{tikzpicture}
\begin{axis}[labmezza, title={Curva valore della capienza},
  xlabel={capienza $k$ (posti)}, ylabel={profitto ottimo (\euro)}]
\addplot+ [mark=*, mark size=1.4pt] table [col sep=comma, x=K, y=profitto]
  {\figdat{cap07_capienza}};
\addplot [grigiomedio, dashed] coordinates {(550, 36000) (550, 60500)};
\node [font=\scriptsize, anchor=east, text=grigiomedio] at (axis cs:540,42000)
  {oltre $q^\circ = 550$ la capienza};
\node [font=\scriptsize, anchor=east, text=grigiomedio] at (axis cs:540,39200)
  {non vale più nulla};
\end{axis}
\end{tikzpicture}
\captionof{figure}{A sinistra: profitto in funzione del prezzo per la domanda
lineare, senza vincolo (parabola grigia tratteggiata) e con capienza $k = 400$
(linea teal): il vincolo taglia la parabola e sposta l'ottimo a $p^* = 160$~\euro.
A destra: profitto ottimo al crescere della capienza; oltre $q^\circ = 550$ posti la
curva diventa piatta e la capacità aggiuntiva non vale più nulla.}
\end{center}

\begin{center}
\begin{tikzpicture}
\begin{axis}[lab, title={Tre funzioni di domanda a confronto},
  xlabel={prezzo (\euro)}, ylabel={domanda attesa (biglietti)},
  ymin=0, ymax=1400]
\addplot [teal, thick] table [col sep=comma, x=p, y=lineare] {\figdat{cap07_domande}};
\addplot [arancio, thick] table [col sep=comma, x=p, y=elast] {\figdat{cap07_domande}};
\addplot [verde, thick] table [col sep=comma, x=p, y=logistica] {\figdat{cap07_domande}};
\addplot [grigiomedio, densely dotted, thick] coordinates {(20, 400) (240, 400)};
\legend{lineare $a - bp$, elasticità costante $\theta p^{-\varepsilon}$,
        logistica, capienza $k = 400$}
\end{axis}
\end{tikzpicture}
\captionof{figure}{Le tre funzioni di domanda del caso di studio a confronto
(biglietti attesi in funzione del prezzo) con la capienza $k = 400$ (linea
punteggiata): la lineare taglia l'asse, l'elasticità costante ha coda lunga, la
logistica descrive la transizione graduale tra domanda alta e bassa.}
\end{center}

\begin{spiegazione}[L'ottimo nel punto di spigolo]
Con elasticità costante l'ottimo \emph{non vincolato} sarebbe
$p = c\,\varepsilon/(\varepsilon - 1) = 36{,}67$~\euro, ma a quel prezzo la domanda
supererebbe di gran lunga la capienza. Per $p$ sotto $79{,}11$~\euro{} si vende
comunque $k$ (conviene quindi alzare il prezzo); sopra, il profitto
$(p - c)\,\theta p^{-\varepsilon}$ decresce. L'ottimo $p^* = 79{,}11$ =
$(\theta/k)^{1/\varepsilon}$ sta nel \emph{punto di spigolo} dove $d(p) = k$: non annulla
nessuna derivata. Morale: mai cercare l'ottimo solo tra i punti stazionari quando ci
sono vincoli.
\end{spiegazione}

\section{Versione multiprodotto}

Due categorie di biglietti con sostituzione (alzare il prezzo della platea spinge
parte della domanda in galleria): $d_1(p_1, p_2) = 500 - 2p_1 + 0{,}6\,p_2$,
$d_2(p_1, p_2) = 900 + 0{,}8\,p_1 - 4p_2$, costi $(30, 15)$~\euro, capacità $(150, 300)$ posti.

\begin{lstlisting}[style=output]
platea   : p1* = 234,04 EUR   q1* = 150/150 (piena)
galleria : p2* = 196,81 EUR   q2* = 300/300 (piena)
profitto totale: 85.148,94 EUR
\end{lstlisting}

\begin{insight}
Entrambe le categorie si riempiono, ma i prezzi non sono ``quelli che svuotano''
ciascuna sala presa da sola: il modello sfrutta la sostituzione, tenendo la platea
cara per spingere domanda verso la galleria. Con prodotti sostituti i prezzi vanno
decisi \emph{congiuntamente}: ottimizzarli uno alla volta lascia soldi sul tavolo.
\end{insight}

\begin{attenzione}
Con funzioni di domanda generali il profitto può non essere concavo: un solver locale
può fermarsi su un ottimo locale. Prassi del laboratorio: studiare dominio e derivate
della funzione \emph{prima} di fidarsi del solver, provare più punti iniziali, e
tracciare sempre il grafico del profitto sull'intervallo rilevante.
\end{attenzione}

\section{Esercizi}

\begin{esercizio}[verifica]
Verificare con Gurobi il valore marginale del posto per $k = 300$ (dalla tabella:
100~\euro) e ricavarlo anche con la formula $p^* - c - k/b$.
\end{esercizio}

\begin{esercizio}[costo marginale]
Come cambiano $p^*$, $q^*$ e il valore del posto se il costo marginale sale a
$c = 60$~\euro? Rifare i tre passi dell'esempio a mano.
\end{esercizio}

\begin{esercizio}[prezzo di riferimento]
Aggiungere la penalità reputazionale $-\gamma\,(p - 120)^2$ con $\gamma = 2$:
quale prezzo bilancia profitto e reputazione? Il problema resta concavo?
\end{esercizio}

\begin{esercizio}[elasticità]
Per la domanda a elasticità costante, tracciare $p^*$ in funzione di
$\varepsilon \in [1{,}3;\, 3]$ con $k = 400$: per quali $\varepsilon$ il vincolo di
capienza resta attivo?
\end{esercizio}

\begin{esercizio}[multiprodotto]
Nel modello a due categorie, di quanto aumenta il profitto con 50 posti di platea in
più? E con 50 di galleria? Quale ampliamento consigliereste, a parità di costo?
\end{esercizio}
